% Evans 8.2: printed pp.437--454 / English PDF pp.454--471.
% Chinese: printed pp.310--320 / scan pp.328--338.
\Needspace{55mm}
\section{Existence of Minimizers}
\label{sec:8-2}
\Needspace{45mm}
\subsection{Coercivity and the direct method}
Fix $g\in W^{1,q}(U;\R^m)$ and the affine class
$\mathcal A=g+W^{1,q}_0(U;\R^m)$, $1<q<\infty$.
The direct method separates three tasks: boundedness of minimizing
sequences, closure of the admissible class under the chosen convergence,
and a lower bound for the energy at the limit.
Suppose
\begin{equation}\label{eq:8-coercivity}
 L(P,z,x)\geq\alpha|P|^q-\beta,\qquad\alpha>0,\ \beta\geq0.
\end{equation}
A bounded energy sequence then has bounded $Du_k$ in $L^q$.
Poincar\'e applied to $u_k-g$ bounds $u_k$ in $W^{1,q}$.
Reflexivity gives a weakly convergent subsequence. The affine class is
weakly closed: the continuous trace map has weakly closed kernel, or use
that a norm-closed linear subspace is weakly closed by Hahn--Banach.
The missing issue is weak lower semicontinuity:
\[
 u_k\rightharpoonup u\text{ in }W^{1,q}
 \quad\Longrightarrow\quad I[u]\leq\liminf_kI[u_k].
\]
Strong norm continuity alone does not imply this inequality.
\Needspace{55mm}
\PDEAnnotation{三个条件各做什么}{强制性给出有界极小化序列；反射性给出弱子列；弱闭性保住边界和约束；弱下半连续性才保证极限达到下确界。四步不能互相替代。}

\PDEBodyDiagram{8.1}{The direct method keeps feasibility and energy in the limit.}
{The arrows represent separate arguments. A finite-energy competitor makes the infimum meaningful; coercivity and Poincare give boundedness. Weak closure preserves constraints, while lower semicontinuity supplies the minimizing inequality.}
{ch08-direct-method}

\Needspace{45mm}
\subsection{Convexity in the gradient and weak lower semicontinuity}
\Needspace{55mm}
\begin{theorem}[Lower semicontinuity; source \S8.2.2, Theorem 1]
\label{thm:8-lsc}
Let $L$ be smooth on gradient and value space times $\overline U$, bounded
below, and convex in $P$ for each $(z,x)$. Then $I$ is sequentially weakly
lower semicontinuous on $W^{1,q}(U;\R^m)$.
No convexity in $z$ is required. The scalar case is the source Theorem 1;
the vector version is used in Theorem 5.
\end{theorem}
\begin{expandedproof}[linearizing only the weakly convergent variable]
Add a constant so $L\geq0$. Pass to a subsequence realizing the liminf;
if it is infinite there is nothing to prove. Compact embedding gives
$u_k\to u$ strongly in $L^q$, and a further subsequence converges a.e.
For $\varepsilon>0$, Egoroff gives a set $E_\varepsilon$ where this
convergence is uniform, with $|U\setminus E_\varepsilon|<\varepsilon$.
Intersect it with $\{|u|+|Du|\leq\varepsilon^{-1}\}$ to obtain
$G_\varepsilon$, whose complement has measure tending to zero.
Convexity gives
\[
 \begin{split}
 I[u_k]\geq{}&\int_{G_\varepsilon}L(Du,u_k,x)\\
 &+\int_{G_\varepsilon}L_P(Du,u_k,x):(Du_k-Du).
 \end{split}
\]
On this set all arguments except $Du_k$ lie in a fixed compact range.
The first integral converges uniformly in its integrand. In the second,
$L_P(Du,u_k,x)$ converges uniformly to $L_P(Du,u,x)$.
Its difference pairs with a uniformly bounded $L^q$ sequence and tends
to zero by H\"older; its fixed limit pairs to zero by weak convergence.
Thus $\liminf I[u_k]\geq\int_{G_\varepsilon}L(Du,u,x)$.
To exhaust $U$ without assuming these sets are nested, fix $R<\infty$:
\[
 \int_{G_\varepsilon}L(Du,u,x)
 \geq\int_U\min\{L(Du,u,x),R\}-R|U\setminus G_\varepsilon|.
\]
Let $\varepsilon\downarrow0$, then $R\uparrow\infty$ by monotone
convergence. This also covers an infinite limit energy.
\end{expandedproof}
\Needspace{55mm}
\PDEAnnotation{为什么只对梯度要求凸}{Rellich 给出 $u_k$ 的强收敛，而 $Du_k$ 通常只有弱收敛。支撑超平面把 $Du_k$ 变成线性出现，正好能使用弱收敛；不能假定梯度有逐点收敛子列。}

\Needspace{55mm}
\begin{theorem}[Existence and uniqueness; source Theorems 2--3, 5--6]
\label{thm:8-direct}
Under \eqref{eq:8-coercivity} and Theorem~\ref{thm:8-lsc}, if
$\mathcal A$ contains $w_0$ with $I[w_0]<\infty$, then $I$ has a
finite minimizer in $\mathcal A$.
If $L=L(P,x)$ is independent of $z$ and uniformly convex in $P$, this
minimizer is unique. These statements apply to scalar and vector maps.
\end{theorem}
\begin{expandedproof}[the limit, and the midpoint deficit]
Choose $I[u_k]\downarrow\inf_{\mathcal A}I$, bounded above by
$I[w_0]+1$ and below by $-\beta|U|$. Coercivity and Poincar\'e give a
weakly convergent subsequence, with limit in $\mathcal A$.
Lower semicontinuity gives $I[u]\leq\inf_{\mathcal A}I$; the reverse
inequality follows from feasibility. This proves existence.
For two minimizers $u,v$, uniform convexity with constant $\theta>0$ gives
\[
 I[(u+v)/2]+\frac\theta8\int_U|Du-Dv|^2
 \leq\frac12(I[u]+I[v]).
\]
This follows by applying the supporting quadratic inequality at the
midpoint and adding the two endpoint inequalities; the linear terms
cancel. Since the midpoint is feasible, minimality forces $Du=Dv$.
Poincar\'e for $u-v\in W^{1,q}_0$ then gives $u=v$.
A merely convex dependence on $P$ with arbitrary $z$ dependence need
not give uniqueness: the value variable may introduce several wells.
\end{expandedproof}
The finite-competitor condition is stated explicitly here; otherwise a
formal minimizer with energy $+\infty$ says nothing about the PDE.

\PDEBodyDiagram{8.2}{Gradient oscillations test the need for convexity.}
{The left graph is $W(p)=(p^2-1)^2$. A zero-trace triangular profile on $(0,1)$ can have slopes $+1,-1$ while its amplitude tends to zero. Its energy is zero, whereas the limit profile has energy $W(0)=1$. The right profile is one exact period, before multiplying its horizontal and vertical scales by $1/k$.}
{ch08-oscillation}

\Needspace{45mm}
\subsection{A Sobolev minimizer solves the Euler--Lagrange equation}
Assume the additional growth bounds, with $q'=q/(q-1)$,
\begin{equation}\label{eq:8-growth}
 \begin{split}
 |L(P,z,x)|&\leq C(1+|P|^q+|z|^q),\\
 |L_P|+|L_z|&\leq C(1+|P|^{q-1}+|z|^{q-1}).
 \end{split}
\end{equation}
\Needspace{55mm}
\begin{theorem}[Weak Euler--Lagrange equation; source Theorems 4 and 7]
A minimizer $u\in\mathcal A$ satisfies
\[
 \int_U L_P(Du,u,x):D\phi+L_z(Du,u,x)\cdot\phi=0
 \quad(\phi\in W^{1,q}_0(U;\R^m)).
\]
\end{theorem}
\begin{expandedproof}[a difference quotient under an integrable bound]
For $|s|\leq1$, the fundamental theorem of calculus expresses
$[L(Du+sD\phi,u+s\phi,x)-L(Du,u,x)]/s$ as an average of derivatives
paired with $(D\phi,\phi)$. Its absolute value is bounded by
\[
 C(1+|Du|^{q-1}+|u|^{q-1}+|D\phi|^{q-1}+|\phi|^{q-1})
                 (|D\phi|+|\phi|),
\]
which lies in $L^1$ by H\"older. Dominated convergence proves the first
variation formula for every admissible $\phi$. Both signs of $s$ are
allowed, so minimality makes the derivative zero. The coefficients
$L_P(Du,u,x),L_z(Du,u,x)$ belong to $L^{q'}$, making this a continuous
weak identity. Without the growth bound, differentiability of the energy
on the whole Sobolev class has not been justified.
\end{expandedproof}
\Needspace{55mm}
\PDEAnnotation{解方程与达到极小}{一阶变分为零只给临界点。若 $L$ 对 $(P,z)$ 联合凸，则支撑不等式配合弱方程可推出全局极小；仅对 $P$ 凸而对 $z$ 非凸时，这个逆命题会失败。}

\Needspace{45mm}
\subsection{Weak continuity of minors and polyconvexity}
\Needspace{55mm}
\begin{lemma}[Weak determinant continuity; source \S8.2.4]
\label{lem:8-det-weak}
If $q>n$ and $u_k\rightharpoonup u$ in $W^{1,q}(U;\R^n)$, then
$\det Du_k\rightharpoonup\det Du$ in $L^{q/n}(U)$.
More generally every $r$-minor converges weakly in $L^{q/r}$, $1\leq r\leq n$.
\end{lemma}
\begin{expandedproof}[induction on minor order]
The gradients are bounded in $L^q$ and Morrey plus compact embedding gives
$u_k\to u$ uniformly after extraction. An $r$-minor is bounded in
$L^{q/r}$ by the polynomial estimate $|M_r(Du_k)|\leq C|Du_k|^r$.
Its Laplace expansion along one fixed component can be written in
 divergence form, with that component undifferentiated and its
$(r-1)$-cofactors inside the flux. The other derivative terms cancel by
the same alternating-index argument as Piola.
For a compactly supported smooth scalar test $\psi$, integration by parts
therefore pairs $u_k^\alpha D\psi$, converging uniformly, with
$(r-1)$-minors. By induction the latter converge weakly in
$L^{q/(r-1)}$. Thus the $r$-minors converge in distributions; the base
$r=1$ is weak gradient convergence. Since $q/r>1$, weak compactness
and density of smooth tests identify every weak subsequential limit
with $M_r(Du)$. This identifies the limit of the original sequence as well.
The same divergence identity extends to Sobolev maps by local smoothing
and H\"older, so it has not been used outside its legitimate class.
\end{expandedproof}
\Needspace{55mm}
\PDEAnnotation{非线性行列式的弱连续}{一般弱收敛不能交换非线性映射。行列式的散度结构把一个因子变为强收敛的 $u_k$，剩余余子式逐阶归纳；$q>n$ 还保证 $L^{q/n}$ 反射。}

The source uses the restricted polyconvex form
\[
 L(P,z,x)=F(P,\det P,z,x),\qquad (P,r)\mapsto F(P,r,z,x)
 \text{ convex}.
\]
The standard broader definition permits all minors. We use the displayed
restricted form, not an unmentioned change of definition in higher dimensions.
\Needspace{55mm}
\begin{theorem}[Polyconvex direct method; source Theorems 8--9]
If $q>n$, $L$ is bounded below, and the smooth representation above holds,
then $I$ is weakly lower semicontinuous. If also
\eqref{eq:8-coercivity} holds and a finite competitor exists, there is a minimizer.
\end{theorem}
\begin{expandedproof}[a supporting plane in an enlarged space]
Use Lemma~\ref{lem:8-det-weak} to get simultaneous weak convergence
of $Du_k$ in $L^q$ and $\det Du_k$ in $L^{q/n}$, and uniform convergence
of $u_k$ along a subsequence. On a good set where $|u|+|Du|$ is bounded,
apply the convex supporting inequality to $F$ in the independent variables
$(P,r)$ at $(Du,\det Du)$. The two linear errors pair with the respective
weakly convergent sequences; the coefficients converge uniformly and
are bounded there. The argument of Theorem~\ref{thm:8-lsc}, including
its truncated exhaustion of the domain, proves the inequality.
Coercivity and weak closure finish the direct method.
For example, $L(P)=|P|^4+(\det P)^2$ for $2\times2$ matrices is of this
form with $q=4$. Polyconvexity does not mean convexity of the nonlinear
map $P\mapsto(P,\det P)$.
\end{expandedproof}

\PDEMarginDiagram{8.M}{秩一方向与全矩阵方向}
{对 $2\times2$ 矩阵，$t\mapsto\det(P+t\,a\otimes b)$ 是仿射函数，但 $t\mapsto\det(tI)=t^2$。秩一方向上的条件不等于整个矩阵空间中的凸性。}
{ch08-rank-one-margin}

\Needspace{45mm}
\subsection{Local minimizers embedded in a field of extremals}
Assume a smooth family $u(x,\lambda)$ solves the scalar Euler--Lagrange
equation for each $\lambda$ in an interval containing zero, with
$u(x,0)=u_0(x)$. Suppose $L$ is convex in $p$. For a smooth
$\vartheta:\overline U\to\R$ taking values in that interval, with
$\vartheta=0$ on $\partial U$, put $w(x)=u(x,\vartheta(x))$.
\Needspace{55mm}
\begin{theorem}[Field of extremals; source \S8.2.5, Theorem 10]
For every such $w$, $I[w]\geq I[u_0]$.
If $u_\lambda$ is bounded away from zero near $\lambda=0$, the inverse
function theorem represents all sufficiently $C^0$-close smooth competitors
with the same trace in this form.
\end{theorem}
\begin{expandedproof}[a calibration flux]
Here $D_xu$ denotes differentiation with $\lambda$ fixed. The chain rule
gives $Dw=D_xu+u_\lambda D\vartheta$ at $\lambda=\vartheta(x)$.
Convexity bounds $L(Dw,w,x)$ below by
$L(D_xu,w,x)+u_\lambda L_p(D_xu,w,x)\cdot D\vartheta$.
Define the vector field
\[
 b_i(x)=\int_0^{\vartheta(x)}
 u_\lambda(x,s)L_{p_i}(D_xu(x,s),u(x,s),x)\,ds.
\]
Differentiating the upper endpoint and using the Euler--Lagrange equation
inside the integral gives
\[
 \begin{split}
 \operatorname{div}b={}&u_\lambda L_p\cdot D\vartheta
 +\int_0^{\vartheta}(L_p\cdot D_xu_\lambda+L_z u_\lambda)\,ds\\
 ={}&u_\lambda L_p\cdot D\vartheta
 +L(D_xu(x,\vartheta),w,x)-L(Du_0,u_0,x).
 \end{split}
\]
Thus $I[w]\geq I[u_0]+\int_U\operatorname{div}b=I[u_0]$, since
$b=0$ on the boundary. The last assertion follows by a uniform local
inverse in $\lambda$ on the compact closed domain. Without that
transversality, the theorem compares only competitors of the stated form.
\end{expandedproof}
